\documentclass{article}
\usepackage{amsmath,hyperref}
\title{Exclusive retry after actual release}
\date{}
\begin{document}
\maketitle

\paragraph{Interface and scope.}
The object is a one-shot civil service, either transferring official possession of a recorder through a docking latch or establishing one authorized load pickup. Several candidates have already staged at the cost specified by the parent I1 fixtures. Exactly one may cause the tracked physical settlement. A capable candidate may not have a functioning channel to another, and receivers only know their own readiness unless the model explicitly supplies a shared observation. Failure of a candidate can leave the exclusive mechanism occupied. The plant can distinguish physical release without a committed result from completion with a result. The recipient needs an effect-side observation or a law guaranteeing both that distinction and when it is actionable. The individual experiments explicitly vary what the receiver can observe or enforce, rather than silently importing these facts.

\paragraph{The single-cut obstruction and public advice.}
Without a shared interlock or new information at the cut, actors with only their own ready/not-ready flag have fixed individual decisions conditional on any fixed common prior state. These define a set $A$ of actors that would start when ready. Exact exclusive commitment is possible if and only if $|A\cap L|=1$ for every admitted ready set $L$. This works when exactly one of two possible actors is ready, but adding the both-ready world makes it impossible if either singleton is still admitted. The proof is not a radio protocol. It assumes the command itself is the onset of the protected effect and there is no feedback or arbitrator within that cut.

For the family of all subsets of size at least $k$ among $n$ candidates, an effect-safe randomized choice must use an $A$ of cardinality at most one almost surely because the all-ready case is admitted. Among the $k$-sets, averaging shows some yield service probability at most $k/n$; a uniform single actor attains it. This is a bounded service probability, not a hard guarantee. For a different design, allow an informed and authorized source to broadcast one reliable public symbol, with each symbol naming a fixed mask $A$. All actors still see only their own readiness plus the symbol. The smallest alphabet has exactly $n-k+1$ symbols. An upper bound uses that many singleton masks, which intersect every live set in the family among them. For the lower bound repeatedly ask the family for a mask that selects exactly one of the remaining actors, starting with all $n$. Delete the selected actor; every earlier witness mask is disjoint from what remains and cannot be reused. There are $n-k+1$ iterations. The cost claim is only $\lceil\log_2(n-k+1)\rceil$ reliable payload bits given a completely informed producer, a shared codebook and effect-time recognition. It does not provide the producer's information, validate physical recipients or cover extra source metadata, lost messages, asynchronous different deliveries, or exactly $k$ ready actors as a different family. This is a scoped instance of the \href{https://citeseerx.ist.psu.edu/document?doi=1c4af6621d5ee28de304dca1bb8f27dd84a35de5&repid=rep1&type=pdf}{ad hoc selective-family problem}, not a priority claim about radio selectors.

\paragraph{The physically enforced alternative.}
A simulated local latch permits competing requests but only one current effect-side claimant. A request must match scope, generation and actual readiness. Once it is held, a new grant or cancel request does not move it. Retirement with no effect and completion of an actual effect have separate local witnesses; both can clear occupancy, but the latter permanently prevents retry of the same one-shot obligation. If a genuinely completed attempt and a failed one look identical at the next cut, retrying duplicates the result in one, while not retrying misses it in the other. Even an authentic indication of bare clearance thus does not establish that retry is lawful. A matching current witness of release without completion does. A wrong scope or generation fails its condition. The explicit latch trace is a truthful physical-and-record model, not an implementation of secure sensors or a theorem that leases are unsafe. \href{https://www.usenix.org/legacy/event/osdi06/tech/full_papers/burrows/burrows_html/index.html}{Chubby's sequencers} already distinguish a lock from effect-side checks on its current generation; the additional consumed datum here is actual physical disposition of the interrupted one-shot effect. A central interlock is a shared physical dependency, not evidence of a decentralized selector.

\paragraph{Exact serial retry when timing is uninformative.}
Fix a healthy available old supervisor and interlock throughout the operating envelope. All candidates are eligible at the first cut, once each; at most $f<n$ attempts fail, with every fixed set of up to $f$ failed candidates admitted. An attempt at route $i$ succeeds with actual completion time cost $s_i$, or fails with authorized re-entry time cost $d_i\ge0$ measured from its start. Here $d_i$ includes any retained physical occupation and delay until its outcome and release may be used; it is not a command-response time. Bound-sharpness assumes the failure prefixes and final success can attain the listed times in one history, or the exact fixed-duration source model. Other source restrictions require the later same-history test.
All actions before the first success observe only completed failed attempts, so every candidate policy has one fixed sequence on the all-fail-prefix branch. Waiting cannot help without new evidence, arrivals or changes. Repeating a failed candidate is not allowed. One can truncate after $f+1$ trials. For an order $\pi$ its exact worst time is
\[
  \max_{j\le f+1}\left(\sum_{i<j}d_{\pi_i}+s_{\pi_j}\right).
\]
For a fixed subset this reduces to \href{https://people.csail.mit.edu/karger/Papers/scheduling.pdf}{the Jackson earliest-due-date ordering} on artificial processing times $d_i$ and tails $s_i-d_i$: sort by decreasing tail. An adjacent exchange proves the order even with negative tails, because shifting a larger tail earlier cannot increase the maximum of the two affected expressions; later prefix sums stay unchanged. For a candidate deadline $D$, scan that order and retain the least cumulative failed time of any feasible prefix with each selected cardinality. Adding an item requires prefix time plus its $s_i$ to be at most $D$; a larger prefix with the same count cannot be better. This gives exact cardinality-$f+1$ selection in $O(n(f+1))$ after sorting and an integer optimal deadline by binary search. Zero-failure and tie cases use the same proof without claiming a better published optimizer.

With two routes $(s,d)=(4,8),(5,2)$, the actual best order is second then first, guaranteeing six ticks. Replacing the first failed hold's eight with a one-tick cancel response makes the first route look preferable. In the failed-first source history the real continuation needs thirteen ticks, or attempts an inadmissibly early re-entry if the logical time is used directly. A genuine latch refuses that early attempt; safety in that case does not imply the promised service deadline. The same earlier duty can stay live only if the old supplier's physical cap also covers the resulting schedule.

\paragraph{Retirement itself can reveal what is still feasible.}
A different fixed finite source law permits correlated route failures and different times of physical retirement. Each attempt has a noninterruptible outcome; no new authority follows bare silence. At a visible state, the controller acts on a cell of possible source rows. Keep actual elapsed time \emph{indexed by each row} even if the receiver has no trustworthy elapsed-time observation. For each proposed next route evaluate successful physical completion deadlines and split the actual failure rows by permitted retirement reports, not by the globally hidden mode. Choose the minimum worst outcome recursively on untried routes. The proof is induction on that finite set: the first choice is common in its local cell, all source-consistent successor observations must be safe, and each continued chosen policy is at the corresponding cell. Numeric elapsed time need only appear in whichever reports the recipient can actually use. Combining worst prefix from one row with worst next service time from another needlessly loses the paired guarantee (e.g. $2+1$ or $1+2$, not $2+2$). This is a standard finite common-information planning reduction, not a novel construction of general policy synthesis.

For the concrete three-route law, A either fails and verifies release in one tick (only B works), fails and verifies release in two ticks (only C works), or actually commits in one tick. B and C each take one tick if successful and physically hold the resource for five on failure. Timed physical verification gives a three-tick guarantee by trying A and selecting the supported next attempt. An untimed but still authentic release leaves the two failure histories indistinguishable; the optimum becomes seven, including consideration of starting with B or C. The fast claim is a comparison of information regimes, not performance against a direct planner given the same evidence.

A bounded error $\epsilon$ on A's local elapsed-time measurement preserves the three-tick guarantee precisely for $\epsilon<1/2$: the two closed observation ranges remain disjoint. At equality a single common reading 1.5 is possible, and the seven-tick obstruction returns. B's and C's failed holds carry no extra distinctions in this pair of histories. The compiler partitions each finite family of interval observations by endpoints and gaps; in a gap the set of compatible rows is constant, and closed endpoints are checked separately. Each world remains fixed while the measurement error ranges may vary independently within the declared per-attempt supports. It is not a robustness claim under arbitrary hidden correlated clock behavior. The timed value is observed after, not instead of, actual release.

At a fixed known drain rate, a local physical record of drained material can replace elapsed time: one bit distinguishing the two amounts realizes the same cut in this example. The source model derives time from decrementing a retraction distance or discharge residue until a physically safe state. If the drain rates can differ, the same residue of six can clear in one or two ticks, and amount alone no longer recovers the required distinction. An independently certified effective rate restores it in the declared new law. The report's meaning depends on actual effect dynamics; neither an unverified command setting nor adding a clock label to existing logs is such a certification. The observed effect source and decoder are simulated and have not been independently transferred to a different modality in the contract's sense.

\paragraph{Verification and comparisons.}
The pulse/scope model checks 35 randomized local exclusivity regimes and 28 advice cardinalities against LP/MILP; the latter are numerical corroborations of the exact lower-bound proof. There are 56 latch controls, 2,900 fixed failure-order comparisons against full permutation and effect traces, 4,080 correlated-time comparisons against direct deadline feasibility, 680 interval-observation comparisons against expanded immutable readout variants, 1,296 typed source-event checks, and 800 fixed-rate time/material source controls. Lean checks the exchange, early-cancel, one-shot ambiguity and no-common-supplier kernels, not network trust. The best same-information direct checks agree. This is a research result about the information and physical prerequisites of the selector in \href{https://github.com/eternal-november/eos-research-programmes/pull/385}{the two-source prototype}; it does not close R1, I1, the held-out I2, or the moonshot's bounded decentralized recurrence.
\end{document}
